\section{Motivating Example}
\label{sec:1.1}

Monte Carlo methods represent a quantity of interest as a parameter of a distribution and use
a random sample from that distribution to estimate the parameter. As an illustrative example,
this section describes a Monte Carlo method for approximating $\pi$, the area of the unit circle $C$.
Let $(X_1, X_2) \sim \Unif([-1,1] \times [-1,1])$ be a random vector uniformly distributed on the square
$S = [-1,1] \times [-1,1]$. Then,
\[
  p := \Prob\!\left((X_1, X_2) \in C\right)
     = \Expect\!\left[\mathbf{1}_{\{(X_1,X_2)\in C\}}\right]
     = \frac{\text{Area } C}{\text{Area } S} = \frac{\pi}{4}.
\]

Hence, $\pi = 4p$, and we have expressed the quantity of interest $\pi$ in terms of the probability $p$.
Now, we can estimate $p$ ---and hence $\pi$--- using a sample
$(X_1^{(1)}, X_2^{(1)}), \ldots, (X_1^{(N)}, X_2^{(N)})$ of
$N$ independent random vectors uniformly distributed on $S$. Precisely, define
\[
  B_N = \sum_{n=1}^{N} \mathbf{1}_{\left\{(X_1^{(n)}, X_2^{(n)}) \in C\right\}}.
\]

Then, $B_N \sim \operatorname{Binomial}(N,p)$ and we can estimate $p$ by $\widehat{p}_N = B_N/N$ and $\pi$ by
$\widehat{\pi}_N = 4\widehat{p}_N$. For example, as shown in Figure~\ref{fig:1.1}, if $N = 10^4$ and we
observe $B_{10000} = 7854$, then
\[
  \widehat{\pi}_{10000} = 4 \times \frac{7854}{10000} = 3.1416.
\]

\begin{figure}[htbp]
\centering
\includegraphics[width=0.75\textwidth]{figures/fig_01_1}
\caption{\textit{Estimation of $\pi$ with sample size $N = 10^4$.} Here, $B_{10000} = 7854$ draws fell within
the unit circle, leading to an estimate $\widehat{\pi}_{10000} = 3.1416$.}
\label{fig:1.1}
\end{figure}

When $N$ is large, the probability ---\textit{before the random sample is generated}--- that our estimate
is far from the true value is small. For $\alpha \in (0,1)$, let $z_{\alpha/2}$ denote the value satisfying
\begin{equation*}
  \Prob(-z_{\alpha/2} < Z < z_{\alpha/2}) = 1 - \alpha, \qquad Z \sim \Normal(0,1).
\end{equation*}

Then,
\[
  \left(\widehat{p}_N - z_{\alpha/2}\sqrt{\frac{\widehat{p}_N(1-\widehat{p}_N)}{N}},\;
         \widehat{p}_N + z_{\alpha/2}\sqrt{\frac{\widehat{p}_N(1-\widehat{p}_N)}{N}}\right)
\]
is an approximate $1-\alpha$ confidence interval for $p$, and therefore
\[
  \left(\widehat{\pi}_N - z_{\alpha/2}\sqrt{\frac{\widehat{\pi}_N(4-\widehat{\pi}_N)}{N}},\;
         \widehat{\pi}_N + z_{\alpha/2}\sqrt{\frac{\widehat{\pi}_N(4-\widehat{\pi}_N)}{N}}\right)
\]
is an approximate $1-\alpha$ confidence interval for $\pi$. In the example shown in Figure~\ref{fig:1.1}
where $N = 10^4$ and $\widehat{\pi}_{10000} = 3.1416$, using the value $z_{0.025} = 1.96$ corresponding
to a 95\% confidence level yields the confidence interval $(3.1094, 3.1738)$.

Let us summarize:
\begin{enumerate}
  \item We expressed the quantity of interest (in our case $\pi$) as an expectation, using that for any
        event $A$, $\Prob(A) = \Expect[\mathbf{1}_A]$. In this textbook, the quantities we seek to
        approximate via Monte Carlo methods will often take the form of expected values.
  \item We constructed an estimator of the expectation based on a random sample.
  \item We used statistical theory to assess the quality of the estimator. The law of large numbers
        guarantees that the sample approximation converges to the quantity of interest.
        Furthermore, the central limit theorem quantifies the speed of this convergence and can be used
        to construct confidence intervals.
\end{enumerate}
